\introChapter{ระบบเลขฐาน (Number Bases)}

\begin{description}[leftmargin=2.5cm, style=nextline]
    \item[รหัสวิชา] 4091606
    \item[ชื่อวิชา] คณิตศาสตร์สำหรับคอมพิวเตอร์
    \item[หน่วยกิต] 3 (3-0-6)
    \item[บทที่] 4
    \item[ชื่อบท] ระบบเลขฐาน (Number Bases)
    \item[จำนวนชั่วโมง] 9 ชั่วโมง
\end{description}

\vspace{0.3em}
\noindent วัตถุประสงค์การเรียนรู้ประจำบท\\[0.3em]
\noindent เมื่อนักศึกษาเรียนบทนี้จบแล้ว นักศึกษาจะสามารถ
\begin{enumerate}[topsep=0.3em, itemsep=0.25em]
    \item แปลงตัวเลขระหว่างระบบเลขฐาน 2, 8, 10 และ 16 ได้อย่างถูกต้อง
    \item ทำการบวก ลบ คูณ หาร ตัวเลขในระบบเลขฐาน 2 ได้
    \item แปลงจำนวนทศนิยมระหว่างระบบเลขฐานและอธิบายความคลาดเคลื่อนจากการปัดเศษได้
    \item อธิบายการแทนข้อมูลตัวอักษรด้วยรหัส ASCII และ Unicode ได้
    \item ประยุกต์ใช้ความรู้เรื่องระบบเลขฐานในการทำงานกับข้อมูลคอมพิวเตอร์ได้
\end{enumerate}

\vspace{0.3em}
\noindent หัวข้อที่สอน\\[0.3em]
\begin{enumerate}[topsep=0.3em, itemsep=0.25em]
    \item ระบบเลขฐาน 10, 2, 8, 16
    \item การแปลงฐาน
    \item เลขคณิตฐานสอง
    \item การแทนข้อมูลในคอมพิวเตอร์ด้วยรหัสอักขระ
    \item การแปลงจำนวนทศนิยมระหว่างฐาน
    \item การแทนข้อมูลด้วยรหัส ASCII และ Unicode
\end{enumerate}

\vspace{1em}
\noindent สื่อการสอน
\begin{enumerate}[topsep=0.3em, itemsep=0.25em]
    \item สไลด์ประกอบการสอน
    \item ตัวอย่างการคำนวณเลขฐาน
    \item แบบฝึกหัดท้ายบท
\end{enumerate}
